\documentclass[12pt]{article}
\usepackage[utf8]{inputenc}
\usepackage[T1]{fontenc}
\usepackage[margin=1in]{geometry}
\usepackage{setspace}
\setlength{\parindent}{2em}
\setlength{\parskip}{1.25em}
\renewcommand{\baselinestretch}{1.0}
\makeatletter
\newcommand{\@subtitle}{}
\newcommand{\subtitle}[1]{\renewcommand{\@subtitle}{#1}}
\renewcommand{\@maketitle}{%
  \newpage
  \null
  \vskip 1em%
  \begin{center}%
    {\LARGE \@title \par}%
    \ifx\@subtitle\@empty\else
      \vskip 0.5em%
      {\large \@subtitle \par}%
    \fi
  \end{center}%
  \par
  \vskip 1em}
\makeatother
\usepackage{xcolor}
\usepackage[colorlinks=true,urlcolor=blue!55!black,linkcolor=black]{hyperref}

\title{\textit{California v. Texas}, 593 U.S. 659 (2021)}
\subtitle{Case Brief}
\author{}
\date{}

\begin{document}
\maketitle

\section{Memory Jogger}

In \href{https://www.courtlistener.com/opinion/4892571/california-v-texas/}{\textit{California v. Texas}}, the Supreme Court dismissed the third major constitutional challenge to the Affordable Care Act (ACA) by a 7--2 vote, holding that neither individual nor state plaintiffs had Article III standing to challenge an individual health-insurance mandate whose noncompliance penalty Congress had reduced to \$0.

\section{Detailed Case Facts}

When Congress enacted the Patient Protection and Affordable Care Act (ACA), Pub. L. 111--148, 124 Stat. 119, in 2010, it included a minimum essential coverage provision---widely known as the ``individual mandate''---codified at 26 U.S.C. \S 5000A(a). The text instructed applicable individuals to maintain qualifying health coverage, while subsections (b) and (c) originally imposed an income-scaled ``shared responsibility payment'' on non-exempt taxpayers who failed to comply, collected through federal income tax returns (\href{https://www.law.cornell.edu/supremecourt/text/19-840}{Cornell LII}).

Two earlier Supreme Court decisions shaped the statutory backdrop. First, in \textit{National Federation of Independent Business v. Sebelius} (\textit{NFIB}), 567 U.S. 519 (2012), five Justices rejected Congress's authority under the Commerce Clause and Necessary and Proper Clause to compel the purchase of health insurance, yet a 5--4 majority led by Chief Justice Roberts upheld \S 5000A under Congress's Article I, \S 8, cl. 1 Taxing Power because the exaction raised federal revenue. Second, in \textit{King v. Burwell}, 576 U.S. 473 (2015), the Court preserved ACA premium tax credits on federally facilitated insurance exchanges (\href{https://www.law.cornell.edu/supremecourt/text/19-840}{Cornell LII}).

In 2017, after broader legislative efforts to repeal the ACA failed in the Senate, Congress passed the Tax Cuts and Jobs Act (TCJA), Pub. L. 115--97, \S 11081, 131 Stat. 2092. Effective beginning tax year 2019, the TCJA amended 26 U.S.C. \S 5000A(c) to set the shared responsibility payment at 0\% of income and \$0, while leaving the coverage language of \S 5000A(a) in the U.S. Code. Following that change, the Internal Revenue Service (IRS) removed the health-coverage check-box from Form 1040 and discontinued Form 8965 (\href{https://www.law.cornell.edu/supremecourt/text/19-840}{Cornell LII}).

In February 2018, Texas and 17 other States (Alabama, Arizona, Arkansas, Florida, Georgia, Indiana, Kansas, Louisiana, Mississippi via its Governor, Missouri, Nebraska, North Dakota, South Carolina, South Dakota, Tennessee, Utah, and West Virginia; Wisconsin and Maine initially participated in the district court before dismissal) sued the United States, the Department of Health and Human Services (HHS), the IRS, and federal officials. Two Texas residents, Neill Hurley and John Nantz, soon joined as individual plaintiffs. The plaintiffs argued that once the penalty dropped to \$0, \S 5000A(a) ceased to raise revenue, could no longer be sustained under the Taxing Power, was therefore unconstitutional, and was inseverable from the rest of the ACA. Because the Trump Administration declined to defend the mandate's constitutionality, California, 15 other States, and the District of Columbia intervened to defend the statute, later joined on appeal by four additional States and the U.S. House of Representatives (\href{https://www.law.cornell.edu/supremecourt/text/19-840}{Cornell LII}; \href{https://harvardlawreview.org/print/vol-135/california-v-texas/}{Harvard Law Review}).

\section{Procedural History}

\begin{itemize}
  \item \textbf{U.S. District Court for the Northern District of Texas:} Treating plaintiffs' filing as a motion for partial summary judgment in \textit{Texas v. United States}, 340 F. Supp. 3d 579 (N.D. Tex. 2018), the District Court ruled that the individual plaintiffs had Article III standing (without deciding state standing), held \S 5000A(a) unconstitutional after the TCJA zeroed out the penalty, and declared the provision inseverable from the entire ACA. The court stayed its declaratory judgment pending appeal, 352 F. Supp. 3d 665 (N.D. Tex. 2018) (\href{https://www.law.cornell.edu/supremecourt/text/19-840}{Cornell LII}).
  \item \textbf{U.S. Court of Appeals for the Fifth Circuit:} In \textit{Texas v. United States}, 945 F.3d 355 (5th Cir. 2019), a 2--1 panel majority (Judge Elrod joined by Judge Engelhardt; Judge King dissenting) held that both the individual and state plaintiffs had standing and affirmed that \S 5000A(a) was unconstitutional. However, the panel concluded that the District Court's wholesale inseverability analysis was incomplete and remanded for a careful, provision-by-provision severability inquiry. The Fifth Circuit denied rehearing en banc, 945 F.3d 182 (5th Cir. 2020) (\href{https://harvardlawreview.org/print/vol-135/california-v-texas/}{Harvard Law Review}).
  \item \textbf{U.S. Supreme Court:} California and the defending state intervenors petitioned for certiorari (No. 19--840), and Texas and its co-plaintiffs filed a conditional cross-petition (No. 19--1019). The Court granted both petitions on March 2, 2020, heard oral argument on November 10, 2020, and issued its decision on June 17, 2021 (\href{https://www.law.cornell.edu/supct/cert/19-840}{Cornell LII Bulletin}). In February 2021, the incoming Biden Administration informed the Court that the Federal Government now regarded \S 5000A(a) as severable.
  \item \textbf{Disposition:} Reversed the Fifth Circuit's judgment as to standing, vacated the judgment, and remanded with instructions to dismiss (\href{https://www.courtlistener.com/opinion/4892571/california-v-texas/}{CourtListener}).
\end{itemize}

\section{Judicial Votes}

The Court ruled \textbf{7--2} to dismiss for lack of standing (\href{https://www.law.cornell.edu/supremecourt/text/19-840}{Cornell LII}):

\begin{itemize}
  \item \textbf{Majority (7):} Justice Breyer authored the opinion of the Court, joined by Chief Justice Roberts and Justices Thomas, Sotomayor, Kagan, Kavanaugh, and Barrett.
  \item \textbf{Concurrence (1):} Justice Thomas filed a concurring opinion while joining the Court's opinion in full.
  \item \textbf{Dissent (2):} Justice Alito filed a dissenting opinion, joined by Justice Gorsuch.
\end{itemize}

\section{Holding}

\textbf{Judgment:} Neither the individual plaintiffs nor the state plaintiffs established Article III standing to challenge 26 U.S.C. \S 5000A(a) because they failed to show a past or future injury fairly traceable to any federal official's actual or threatened enforcement of that provision. The Court therefore did not reach whether \S 5000A(a) is unconstitutional or severable (\href{https://www.law.cornell.edu/supremecourt/text/19-840}{Cornell LII}).

\textbf{Rule Announced:} Article III standing to challenge a federal statute requires a concrete, particularized injury fairly traceable to the defendant's actual or threatened enforcement of the specific statutory provision attacked as unconstitutional, and likely to be redressed by a judicial remedy operating against specific parties. Textually unenforceable statutory language carrying no penalty or enforcement mechanism cannot by itself establish traceability or redressability (\href{https://www.law.cornell.edu/supremecourt/text/19-840}{Cornell LII}).

\textbf{Holding vs. Dicta:} The binding holding is strictly jurisdictional: (1) individuals who buy health insurance to follow a \$0-penalty mandate suffer no injury traceable to federal enforcement; (2) states alleging indirect enrollment costs in Medicaid, CHIP, or state employee health plans must produce concrete summary-judgment evidence that an unenforceable mandate predicts third-party enrollment choices; and (3) states' direct administrative or compliance expenses imposed by independent ACA provisions (such as 26 U.S.C. \S\S 4980H, 6055, and 6056) are not traceable to \S 5000A(a). Because the majority rejected the ``standing-through-inseverability'' theory solely on forfeiture grounds, the Court did not resolve whether that theory could succeed when timely preserved---a limitation explicitly flagged in Justice Thomas's concurrence (\href{https://www.law.cornell.edu/supremecourt/text/19-840}{Cornell LII}).

\section{Analysis of Opinions}

\subsection{Majority Opinion (Justice Breyer)}

Writing for seven Justices, Justice Breyer grounded the decision in Article III, \S 2's ``Cases'' and ``Controversies'' requirement and the established standing elements of \textit{Lujan v. Defenders of Wildlife}, 504 U.S. 555 (1992), and \textit{DaimlerChrysler Corp. v. Cuno}, 547 U.S. 332 (2006): injury in fact, traceability, and redressability (\href{https://www.law.cornell.edu/supremecourt/text/19-840}{Cornell LII}).

\textbf{Individual Plaintiffs:} Even assuming that paying monthly premiums counted as a pocketbook injury, Hurley and Nantz failed the traceability requirement. Because Congress reduced the \S 5000A(c) penalty to \$0 and \S 5000A(g) authorizes IRS enforcement only for nonpayment of a penalty, no federal officer could enforce \S 5000A(a) against them. Earlier ACA challenges arose when the penalty was active and enforceable, whereas here only unenforceable statutory text remained. Redressability underscored the same defect: federal remedies operate against specific officials rather than abstract legal rules, and the Declaratory Judgment Act, 28 U.S.C. \S 2201, does not independently confer Article III jurisdiction. With no enforcement conduct to enjoin, declaring \S 5000A(a) invalid would amount to an advisory opinion (\href{https://www.law.cornell.edu/supremecourt/text/19-840}{Cornell LII}).

\textbf{State Plaintiffs' Indirect Injury Claim:} The States contended that \S 5000A(a) drove additional residents to enroll in state-funded programs like Medicaid, the Children's Health Insurance Program (CHIP), and state-employee health plans. Justice Breyer rejected this theory both as a matter of law (no federal enforcement of \S 5000A(a) caused enrollment) and as a failure of summary-judgment proof. Under \textit{Department of Commerce v. New York}, 588 U.S. 752 (2019), and \textit{Clapper v. Amnesty Int'l USA}, 568 U.S. 398 (2013), standing theories hinging on independent third-party decisions require concrete proof that third parties will react predictably. Because Medicaid and CHIP already provide substantial no-cost or low-cost benefits, logic did not suggest that a toothless mandate would induce enrollment. Only four of the States' 21 declarations tied costs to the mandate, and all four referred to the pre-2019 era when a monetary penalty existed. Likewise, a single predictive line in a 2017 Congressional Budget Office (CBO) report noting that a small number of people might still buy coverage to obey the law failed to demonstrate that those individuals would enroll in state-funded plans rather than private markets (\href{https://www.law.cornell.edu/supremecourt/text/19-840}{Cornell LII}).

\textbf{State Plaintiffs' Direct Injury Claim:} The States also cited direct expenses from IRS reporting forms under 26 U.S.C. \S\S 6055 and 6056, the employer mandate in \S 4980H, coverage for employees working 30--39 hours per week, and continuing Medicaid coverage for former foster youth. Justice Breyer explained that those separate statutory provisions operate independently of \S 5000A(a), borrowing only \S 5000A(f)'s definition of ``minimum essential coverage.'' Because the States did not challenge those other provisions as independently unconstitutional, compliance costs under them were not fairly traceable to \S 5000A(a). Finally, the Court refused to entertain the novel ``standing-through-inseverability'' theory advanced late in the litigation by the Federal Government and the dissent because plaintiffs had not raised it in the lower courts or at the certiorari stage (\href{https://www.law.cornell.edu/supremecourt/text/19-840}{Cornell LII}).

\subsection{Concurring Opinion (Justice Thomas)}

Justice Thomas joined the Court's opinion in full while writing separately to acknowledge the force of Justice Alito's historical critique of \textit{NFIB} and \textit{King v. Burwell}. He noted the sharp reversal by the ACA's defenders, who had insisted during \textit{NFIB} that the individual mandate was indispensable to guaranteed-issue and community-rating reforms, yet now portrayed it as an isolated, severable sentence (\href{https://www.law.cornell.edu/supremecourt/text/19-840}{Cornell LII}).

Even so, Justice Thomas concluded that the plaintiffs failed to tie any concrete injury to unlawful federal action. Turning to Justice Alito's ``standing-through-inseverability'' theory, Justice Thomas observed that such a theory might be viable in a future case if framed as a merits question of statutory construction rather than a freestanding judicial remedy. He declined to apply it here for four reasons: (1) plaintiffs did not raise it below; (2) the States did not clearly develop it in their Supreme Court briefs; (3) earlier cases cited by the dissent were unreasoned ``drive-by jurisdictional rulings'' lacking precedential force under \textit{Steel Co. v. Citizens for Better Environment}, 523 U.S. 83 (1998); and (4) the Court's severability doctrine had inconsistently treated inseverability as remedial rather than statutory (\href{https://www.law.cornell.edu/supremecourt/text/19-840}{Cornell LII}).

\subsection{Dissenting Opinion (Justice Alito, Joined by Justice Gorsuch)}

Justice Alito opened by characterizing the decision as the third chapter of an ACA trilogy in which the Court repeatedly rescued the statute through strained legal reasoning. Focusing on the 18 plaintiff States, he argued that they plainly suffered concrete, ongoing pocketbook injuries---including hundreds of thousands of dollars in annual Forms 1094/1095 reporting expenses under \S\S 6055 and 6056 and tens of millions of dollars providing coverage to adult children up to age 26 under 42 U.S.C. \S 300gg--14(a)---backed by steep federal fines (\href{https://www.law.cornell.edu/supremecourt/text/19-840}{Cornell LII}).

\textbf{Traceability and Inseverability Standing:} Justice Alito accused the majority of misreading \textit{Allen v. Wright}, 468 U.S. 737 (1984), by requiring plaintiffs to trace their injury to enforcement of the specific \textit{unconstitutional provision} rather than to the defendant's \textit{challenged conduct} alleged to be unlawful. Invoking 141 years of Supreme Court cases---including \textit{Seila Law LLC v. CFPB} (2020), \textit{Free Enterprise Fund v. PCAOB} (2010), \textit{Minnesota v. Mille Lacs Band of Chippewa Indians} (1999), \textit{New York v. United States} (1992), \textit{Alaska Airlines, Inc. v. Brock} (1987), and the \textit{Trade-Mark Cases} (1879)---he argued that a party burdened by Provision A has standing to contend that Provision A is unenforceable because it is inseverable from an unlawful Provision B. He also disputed forfeiture, citing the States' amended complaint and appellate briefs, and pointed out in footnote 9 that a jurisdictional dismissal does not preclude a renewed lawsuit (\href{https://www.law.cornell.edu/supremecourt/text/19-840}{Cornell LII}).

\textbf{Merits and Severability:} Reaching the merits, Justice Alito argued that a \$0 penalty can never generate revenue---the defining feature of a tax under \textit{NFIB}---and that the mandatory verb ``shall'' precluded treating \S 5000A(a) as merely precatory. On severability, he applied both the \textit{NFIB} joint dissent's two-step test and textualist statutory-interpretation principles, pointing to Congress's express finding in 42 U.S.C. \S 18091(2)(I) that the mandate was ``essential'' to the ACA's insurance reforms. Concluding that the 2017 TCJA did not demonstrate congressional intent to keep the rest of the ACA intact, he would have barred federal officials from enforcing the state-burdening ACA provisions against the plaintiff States (\href{https://www.law.cornell.edu/supremecourt/text/19-840}{Cornell LII}).

\section{Examples: Future Applications}

\begin{enumerate}
  \item \textbf{Same-Side Application (Zero-Penalty Federal Composting Directive):} Congress enacts a statute stating that every household ``shall'' separate food waste for composting, but sets the civil penalty at \$0 and grants no agency authority to inspect or fine noncompliant households. A homeowner who buys composting bins out of civic duty sues for a declaratory judgment that Congress exceeded its Commerce Clause power. \textit{Outcome:} \textbf{No standing.} Under \textit{California v. Texas}, voluntary expenditures prompted by unenforceable statutory text are not fairly traceable to federal enforcement, and a bare declaratory judgment would be an advisory opinion.
  \item \textbf{Same-Side Application (Speculative State Costs From an Unenforceable Safety Mandate):} A federal statute instructs homeowners to install smart smoke detectors without attaching any fine, tax, or enforcement mechanism. Several States sue, alleging higher outlays under their state-run home-safety rebate programs because residents will claim rebates to obey the federal statute. \textit{Outcome:} \textbf{No standing.} Without concrete proof that an unenforceable federal directive---rather than the financial attractiveness of the state rebate itself---drives third-party participation, the causal chain is too speculative under \textit{California v. Texas} and \textit{Clapper}.
  \item \textbf{Opposite-Side Application (Enforceable Monetary Assessment):} Congress amends the composting statute in Example 1 to impose a \$200 annual IRS assessment on noncompliant households, enforceable by offsetting federal tax refunds. A homeowner who buys composting gear to avoid the charge sues to enjoin IRS collection. \textit{Outcome:} \textbf{Standing exists.} Because a federal official possesses concrete enforcement authority creating a realistic threat of financial sanction, the injury is traceable to threatened government action and redressable by an injunction.
  \item \textbf{Opposite-Side Application (Collateral Federal Sanctions Despite a \$0 Fine):} Congress reduces the direct monetary fine for violating a federal cybersecurity certification rule to \$0, but a companion clause automatically bars any noncompliant vendor from bidding on federal procurement contracts. A contractor that spends \$50,000 to maintain certification sues the procurement agency. \textit{Outcome:} \textbf{Standing exists.} Even with a \$0 fine, federal debarment supplies a concrete enforcement mechanism that a court can enjoin.
  \item \textbf{Fence-Sitter (Timely Preserved Standing-Through-Inseverability With an Express Clause):} Congress passes an energy omnibus act containing Provision A (an enforceable \$25,000-per-day reporting rule on utilities), Provision B (an unenforceable \$0 directive telling drivers to buy electric vehicles), and an explicit inseverability clause stating that Provision A is void if Provision B is held unconstitutional. A utility burdened by Provision A sues the Department of Energy, squarely pleading and briefing standing-through-inseverability from the start of the case. \textit{Outcome:} \textbf{Uncertain / Fence-sitter.} Because both the majority and Justice Thomas left properly preserved standing-through-inseverability undecided on forfeiture grounds, a court would have to decide whether an express statutory link transforms inseverability into a cognizable merits-based standing theory.
\end{enumerate}

\section{Critique}

\textbf{Scholarly Criticism and Doctrinal Ambiguity:} Legal scholars across ideological lines observed that \textit{California v. Texas} defused a high-stakes political controversy through Article III gatekeeping while leaving state-standing doctrine unsettled. In its Leading Case comment, the \href{https://harvardlawreview.org/print/vol-135/california-v-texas/}{\textit{Harvard Law Review}} (135 Harv. L. Rev. 343 (2021)) highlighted that neither the majority, the dissent, the lower courts, nor the litigants engaged with \textit{Massachusetts v. EPA}, 549 U.S. 497 (2007), and its ``special solicitude'' standard for state plaintiffs. Leaving that precedent unmentioned---neither applying nor overruling it---created practical uncertainty for lower courts forced to guess whether ``special solicitude'' had been functionally abandoned through stealth overruling (\href{https://harvardlawreview.org/print/vol-135/california-v-texas/}{Harvard Law Review}).

\textbf{Logical Weaknesses:} Critics also debated the analytical boundary between standing and the merits. Justice Alito's objection---that requiring traceability to enforcement of the \textit{specific unconstitutional provision} blurs threshold jurisdiction with substantive severability---exposes a tension with earlier cases where the Court entertained inseverability challenges without questioning standing. Conversely, defenders of the majority argued that permitting any litigant regulated by a valid provision to manufacture standing by attacking an unrelated, unenforceable clause elsewhere in an omnibus statute would invite sweeping judicial oversight of symbolic legislation.

\textbf{Progressive Perspective:} Progressive commentators welcomed the 7--2 dismissal as a sensible institutional outcome that safeguarded health insurance for tens of millions of Americans and prevented courts from dismantling a landmark statute after Congress itself declined to repeal it in 2017. At the same time, progressive scholars noted that the decision formed part of a broader October Term 2020 tightening of Article III standing (alongside \textit{TransUnion LLC v. Ramirez} and \textit{Carney v. Adams}) that could later impede consumers, environmental groups, and state attorneys general challenging federal deregulation (\href{https://harvardlawreview.org/print/vol-135/california-v-texas/}{Harvard Law Review}).

\textbf{Originalist and Textualist Perspective:} Originalists and textualists split along the same fault line visible between Justices Thomas, Kavanaugh, and Barrett in the majority and Justices Alito and Gorsuch in dissent. Judicial-restraint textualists applauded the Court for honoring traditional Article III limits---courts decide concrete disputes and enjoin unlawful official acts rather than issuing advisory declarations against unenforceable text. Other originalists shared Justice Alito's frustration, contending that the Court applied standing rules selectively to insulate the ACA and left an enumerated-powers violation unredressed once the \$0 penalty erased the Taxing Power rationale of \textit{NFIB}.

\section{Key Quotations}

\begin{quote}
``The provision is unenforceable. There is no one, and nothing, to enjoin.'' --- Justice Breyer, for the Court (\href{https://www.law.cornell.edu/supremecourt/text/19-840}{Cornell LII}).
\end{quote}

\begin{quote}
``To find standing here to attack an unenforceable statutory provision would allow a federal court to issue what would amount to `an advisory opinion without the possibility of any judicial relief.' '' --- Justice Breyer, for the Court (\href{https://www.law.cornell.edu/supremecourt/text/19-840}{Cornell LII}).
\end{quote}

\begin{quote}
``A penalty might have led some inertia-bound individuals to enroll. But without a penalty, what incentive could the provision provide?'' --- Justice Breyer, for the Court (\href{https://www.law.cornell.edu/supremecourt/text/19-840}{Cornell LII}).
\end{quote}

\begin{quote}
``Although this Court has erred twice before in cases involving the Affordable Care Act, it does not err today.'' --- Justice Thomas, concurring (\href{https://www.law.cornell.edu/supremecourt/text/19-840}{Cornell LII}).
\end{quote}

\begin{quote}
``A penalty is a tax. The United States is a State. And 18 States who bear costly burdens under the ACA cannot even get a foot in the door to raise a constitutional challenge.'' --- Justice Alito, dissenting (\href{https://www.law.cornell.edu/supremecourt/text/19-840}{Cornell LII}).
\end{quote}

\end{document}
